\section{Return levels and reconstructed coefficients together}
Implement \texttt{h263\_quantize8}. Pass a mapped eight-byte-aligned 32-byte
request in \texttt{\$a0} and a mapped sixteen-byte-aligned 256-byte result in
\texttt{\$a1}. The request contains unsigned native-endian doublewords:
\begin{center}
\tagpdfsetup{table/header-rows={1}}
\begin{tabular}{rll}
\textbf{Offset} & \textbf{Field} & \textbf{Meaning}\\
0 & Coefficient pointer & Two-byte-aligned, 128 active bytes\\
8 & Quantizer parameter & 1 through 31\\
16 & INTRA selector & 0 INTER; 1 INTRA\\
24 & Arithmetic mode & 0 scalar; 1 MSA reconstruction
\end{tabular}
\end{center}
The input contains 64 signed halfwords in natural row-major coefficient order.
Every value must be in $[-2048,2047]$. Input, request and result are disjoint.
Their complete extents must be mapped and must not wrap the address space.
Supply a sixteen-byte-aligned stack and valid private space. Preserve saved
general registers, the incoming stack pointer and every input/request byte.
Vector registers are caller-saved.

Validate record alignment and extents, input alignment and extent, and every
full-width control before reading coefficient samples. Invalid metadata returns
\texttt{\$v0}=1 and \texttt{\$v1}=0, retaining the complete result and avoiding
sample reads. After valid metadata, inspect the signed coefficient domain. An
invalid sample returns status two and count zero, also retaining the result.
Reads needed to discover that invalid sample are permitted.

On success, the first 128 result bytes hold 64 signed halfword levels. For INTRA,
the first halfword instead holds the unsigned DC code as a positive number from
1 through 255, excluding 128. The next 128 bytes hold 64 signed reconstructed
coefficients ready for the inverse transform. Return status zero and a count of
nonzero ordinary levels, excluding INTRA DC. The count can therefore be zero
for an INTRA block whose reconstructed DC is nonzero. It is a coefficient count,
not a coded-block-pattern mask or a promise about the final number of bits.

Mode zero uses scalar arithmetic throughout. Mode one retains scalar level
selection and uses MSA lanes for ordinary inverse quantization. Each mode obeys
the same exact exercise policy. Behavioral grading checks outputs and state;
matching outputs alone does not demonstrate that a submission used MSA. Inspect
that instructional objective separately when assessing the vector implementation.

Start with \path{construction/h263-quantization-construction/main.asm}. Prepare another
observation while preserving your instructions and coefficient data:
\begin{verbatim}
python3 construction/h263-quantization-construction/prepare.py \
    --source my-quantizer.asm --quant 8 --intra 1 \
    --mode 1 --output quantizer-q8.asm
\end{verbatim}
The helper changes the typed request controls and refuses an existing output.
The provided caller aligns its stack, installs the input pointer and calls your
entry. At \texttt{quant\_caller\_return}, the result is complete and return
registers still contain status and count. Named status/count fields are copied
by subsequent instructions. The completed default prints 255 and exits zero;
the unfinished entry reports development status four and exits one.

A successful printed first value checks only one code. Use the public checker to
inspect every level and reconstructed coefficient, as well as preserved memory
and registers. Keep natural order here; scan order and bit emission come later.
